\documentclass[a4paper]{article}

\usepackage[fontset=fandol]{ctex}
\usepackage{amsmath, amssymb}
\usepackage{graphicx}
\usepackage[margin=2.5cm]{geometry}
\usepackage[most]{tcolorbox}
\usepackage{listings}
\usepackage{hyperref}
\usepackage{booktabs}
\usepackage{float}
\usepackage{tikz}

\setlength{\parindent}{2em}
\setlength{\parskip}{0.35em}
\sloppy

\newtcolorbox{knowledgebox}[1]{
    enhanced, colback=blue!5!white, colframe=blue!75!black, colbacktitle=blue!75!black,
    coltitle=white, fonttitle=\bfseries, title={#1}, attach boxed title to top left={yshift=-2mm, xshift=2mm},
    boxrule=1pt, sharp corners
}

\newtcolorbox{importantbox}[1]{
    enhanced, colback=yellow!10!white, colframe=yellow!80!black, colbacktitle=yellow!80!black,
    coltitle=black, fonttitle=\bfseries, title={#1}, sharp corners
}

\newtcolorbox{warningbox}[1]{
    enhanced, colback=red!5!white, colframe=red!75!black, colbacktitle=red!75!black,
    coltitle=white, fonttitle=\bfseries, title={#1}, sharp corners
}

\lstset{
    basicstyle=\ttfamily\small,
    keywordstyle=\color{blue},
    stringstyle=\color{red!60!black},
    commentstyle=\color{green!60!black},
    breaklines=true,
    frame=single,
    numbers=left,
    numberstyle=\tiny\color{gray},
    captionpos=b,
    extendedchars=false
}

\newcommand{\notetitle}{批次标准化 Batch Normalization}
\newcommand{\notesubtitle}{李宏毅《机器学习》2025 第五节：让训练地形更好走}
\newcommand{\noteauthors}{基于李宏毅老师课程整理}
\newcommand{\notedate}{\today}
\newcommand{\videochannel}{李宏毅机器学习深度学习课程（Bilibili）}
\newcommand{\videoduration}{约 31 分钟}
\newcommand{\videourl}{https://www.bilibili.com/video/BV1TAtwzTE1S?p=14}

\begin{document}
\pagestyle{empty}

\begin{center}
    \vspace*{1.5cm}
    {\Huge\bfseries \notetitle\par}
    \vspace{0.4cm}
    {\LARGE \notesubtitle\par}
    \vspace{0.8cm}
    \includegraphics[width=0.62\textwidth]{video.jpg}\par
    \vspace{0.8cm}
    {\large \noteauthors\par}
    {\large \notedate\par}
    \vspace{0.8cm}
    \begin{tcolorbox}[width=0.82\textwidth,center,sharp corners,
                      colback=black!3!white,colframe=black!50!white]
        \centering
        \textbf{视频信息}\\[3pt]
        频道：\videochannel \\
        时长：\videoduration \\
        链接：\href{\videourl}{\videourl}
    \end{tcolorbox}
\end{center}

\newpage
\tableofcontents
\newpage
\pagestyle{plain}

% =====================================================================
\section{为什么 Batch Normalization 出现在这里}
% =====================================================================

这堂课接在“类神经网络训练不起来怎么办”的系列后面。前面几讲讨论过 learning rate、momentum、adaptive learning rate、loss function 等训练技巧，它们大多是在 optimizer 或 objective 上动手。Batch Normalization 则更像是直接改造网络内部的计算流程，让训练时面对的 error surface 更容易优化。

老师用“把山铲平”来描述这个想法：如果 loss landscape 在某些方向很陡、某些方向很平，gradient descent 就容易在峡谷里左右震荡，或者只能用很小的 learning rate 慢慢走。BN 的目标不是换一个 optimizer，而是让中间表示的尺度更稳定，使训练地形相对好走。

\begin{importantbox}{Batch Normalization 的核心问题}
    深度网络每一层都会改变资料分布。即使原始输入已经做过 normalization，经过前几层线性变换和 activation 后，后面层收到的 hidden features 仍可能有不同尺度、不同均值、不同方差。BN 尝试在训练过程中持续把这些中间特征拉回比较统一的尺度。
\end{importantbox}

\subsection{本章小结}

BN 属于“改变网络内部地形”的训练技巧。它不是单纯预处理输入，而是在深层网络的中间层反复做 normalization。

% =====================================================================
\section{尺度差异如何改变 Error Surface}
% =====================================================================

先看最简单的线性模型：
\[
    y = b + w_1x_1 + w_2x_2.
\]
如果 $x_1$ 的值通常很小，例如 1、2；而 $x_2$ 的值通常很大，例如 100、200，那么同样大小的参数改变会造成非常不同的输出变化。改变 $w_1$ 一点点，对 $y$ 的影响可能很小；改变 $w_2$ 一点点，对 $y$ 的影响可能很大。

从梯度角度看，
\[
    \frac{\partial L}{\partial w_i}
    =
    \sum_r
    \frac{\partial L}{\partial y^r}
    x_i^r.
\]
当某个 dimension 的 $x_i$ 尺度特别大，它对应方向的 gradient 和 loss curvature 都可能特别大。于是 error surface 变成长而窄的椭圆：一个方向平滑，另一个方向陡峭。

\begin{figure}[H]
    \centering
    \includegraphics[width=0.86\textwidth]{figures/fig01_feature_scale_landscape.jpg}
    \caption{不同输入维度尺度差异过大时，loss landscape 可能在不同方向呈现不同斜率；把 feature 调到相近范围有助于让地形更圆滑。\protect\footnotemark}
    \label{fig:feature-scale}
\end{figure}
\footnotetext{视频画面时间区间：约 05:20--06:05；字幕对应“input feature 每一个 dimension 的 scale 差距很大”“产生不同方向斜率非常不同的 error surface”。}

\begin{knowledgebox}{为什么椭圆地形难走}
    在狭长谷底中，若 learning rate 太大，参数会在陡峭方向来回震荡；若 learning rate 太小，平缓方向又前进太慢。Feature normalization 不保证 loss 变成完美圆形，但它能减少某些纯粹由输入尺度造成的坏条件数。
\end{knowledgebox}

\subsection{本章小结}

不同 dimension 的尺度差异会被参数更新放大成优化难度。把 feature 调整到相近均值和方差，是改善训练地形的第一步。

% =====================================================================
\section{Feature Normalization}
% =====================================================================

给定训练资料
\[
    x^1,x^2,\ldots,x^R,
\]
其中每个 $x^r$ 都是一个多维向量。对第 $i$ 个 dimension，可以计算训练集中该维度的平均值 $m_i$ 和标准差 $\sigma_i$：
\[
    m_i=\frac{1}{R}\sum_{r=1}^{R}x_i^r,\qquad
    \sigma_i=\sqrt{\frac{1}{R}\sum_{r=1}^{R}(x_i^r-m_i)^2}.
\]
然后把每个样本的该维度改写成
\[
    \tilde x_i^r = \frac{x_i^r-m_i}{\sigma_i}.
\]
这样每个 dimension 的平均值约为 0，方差约为 1。

\begin{figure}[H]
    \centering
    \includegraphics[width=0.86\textwidth]{figures/fig02_feature_normalization_formula.jpg}
    \caption{Feature normalization：对每个维度分别减去平均值、除以标准差，使各维度平均约为 0、方差约为 1。\protect\footnotemark}
    \label{fig:feature-normalization}
\end{figure}
\footnotetext{视频画面时间区间：约 07:50--08:25；字幕对应“减掉这个 dimension 算出来的 mean，再除掉 standard deviation”“平均是零”。}

许多深度学习作业或实作都会先对输入 feature 做 normalization。它有两个实际效果：
\begin{itemize}
    \item 让各个 input dimension 对参数更新的影响不至于仅由尺度决定；
    \item 让模型更容易使用统一的 learning rate，而不是某些参数方向需要极小步长。
\end{itemize}

\begin{warningbox}{Normalization 不是资料泄漏的借口}
    训练时的平均值和标准差应该由训练集估计，再用于 validation 或 test。若用完整资料集、尤其包含测试资料来估计 normalization statistics，就会把测试集分布信息泄漏进训练流程。
\end{warningbox}

\subsection{本章小结}

Feature normalization 是 BN 的前置直觉：每个维度先对齐尺度，优化会更容易。但深度网络的问题不止在输入层。

% =====================================================================
\section{深层网络中为什么还需要 Normalization}
% =====================================================================

即使原始输入 $\tilde x$ 已经 normalization，经过第一层线性变换
\[
    z = W^1\tilde x
\]
之后，$z$ 的各个维度仍可能有不同范围。再经过 activation 得到 $a$，传给下一层 $W^2$ 时，下一层看到的其实是新的 input。换句话说，每一层的 hidden representation 都可能重新产生尺度问题。

\begin{figure}[H]
    \centering
    \includegraphics[width=0.86\textwidth]{figures/fig03_hidden_feature_ranges.jpg}
    \caption{原始输入做过 normalization 不代表 hidden features 也保持同样尺度；经过 $W^1$ 与 activation 后，不同维度仍可能有不同范围。\protect\footnotemark}
    \label{fig:hidden-ranges}
\end{figure}
\footnotetext{视频画面时间区间：约 10:45--11:25；字幕对应“对 Z 做 feature normalization 或对 A 做 feature normalization 都可以”“sigmoid 在零附近斜率比较大”。}

于是可以把 normalization 放在 hidden layer 中。常见选择是对 activation function 前的 pre-activation $z$ 做 normalization：
\[
    \tilde z = \frac{z-\mu}{\sigma}.
\]
之后再送进 activation function。若 activation 是 sigmoid，把 $z$ 拉回 0 附近还有一个额外直觉：sigmoid 在 0 附近斜率较大，比较不容易进入饱和区。

\begin{knowledgebox}{BN 放在 activation 前还是后}
    课堂中说，实作上二者差异未必总是巨大；若使用 sigmoid，通常更推荐放在 activation 前，因为可以让输入 activation 的值靠近斜率较大的区域。现代 CNN 中常见顺序是 convolution、batch normalization、activation。
\end{knowledgebox}

\subsection{本章小结}

深层网络每一层都会制造新的 feature distribution，因此只 normalize 原始输入不够。BN 把 normalization 嵌入网络内部，使 hidden features 也保持较稳定的尺度。

% =====================================================================
\section{Batch Normalization 的计算图}
% =====================================================================

BN 的关键是：hidden feature 的平均值和标准差不是对单一样本计算，而是对一个 mini-batch 计算。设某一层对一个 batch 产生 pre-activation
\[
    z^1,z^2,\ldots,z^m,
\]
其中 $m$ 是 batch size。BN 先计算 batch mean 和 batch variance：
\[
    \mu_B = \frac{1}{m}\sum_{r=1}^{m}z^r,
\]
\[
    \sigma_B^2
    =
    \frac{1}{m}\sum_{r=1}^{m}(z^r-\mu_B)^2.
\]
然后对 batch 中每个样本做 normalization：
\[
    \tilde z^r
    =
    \frac{z^r-\mu_B}{\sqrt{\sigma_B^2+\epsilon}}.
\]
这里 $\epsilon$ 是很小的数，用来避免除以零。

\begin{figure}[H]
    \centering
    \includegraphics[width=0.86\textwidth]{figures/fig04_batch_norm_large_network.jpg}
    \caption{在 deep network 中，$\mu$ 和 $\sigma$ 由一整个 batch 的 hidden features 决定；BN 因此成为网络计算图的一部分。\protect\footnotemark}
    \label{fig:bn-graph}
\end{figure}
\footnotetext{视频画面时间区间：约 17:20--17:55；字幕对应“你一定要有一个够大的 batch 才算得出 mu 跟 sigma”。}

这个图很重要，因为它说明 BN 不是对每笔资料独立做一个固定变换。某个样本的 normalized value 会受同 batch 中其他样本影响；反向传播时，梯度也会通过 $\mu_B$ 和 $\sigma_B$ 连接 batch 内样本。

\begin{warningbox}{Batch size 太小会影响 BN}
    如果 batch size 太小，batch statistics 会很不稳定；极端情况下 batch size 为 1，方差几乎无法提供有意义的尺度信息。小 batch 场景常需要改用 Layer Normalization、Group Normalization，或使用同步 BN 等技巧。
\end{warningbox}

\subsection{本章小结}

BN 的 “batch” 不是名字装饰，而是机制核心：统计量来自 mini-batch，因此训练时样本之间在 normalization 层发生耦合。

% =====================================================================
\section{为什么还要 Gamma 与 Beta}
% =====================================================================

如果只做
\[
    \tilde z^r
    =
    \frac{z^r-\mu_B}{\sqrt{\sigma_B^2+\epsilon}},
\]
每个 dimension 都会被强行变成接近零均值、单位方差。这可能过度限制模型。也许某一层就是需要非零均值，或者需要更大的尺度，才能表示最适合任务的函数。

因此 BN 通常再加入两个可学习参数：scale $\gamma$ 与 shift $\beta$：
\[
    \hat z^r = \gamma \odot \tilde z^r + \beta.
\]
其中 $\odot$ 表示 element-wise multiplication。$\gamma$ 和 $\beta$ 的维度与 feature dimension 相同，训练过程中由 gradient descent 学出。

\begin{figure}[H]
    \centering
    \includegraphics[width=0.86\textwidth]{figures/fig05_gamma_beta_restore.jpg}
    \caption{BN 先标准化，再用可学习的 $\gamma$ 与 $\beta$ 恢复模型需要的尺度和偏移。\protect\footnotemark}
    \label{fig:gamma-beta}
\end{figure}
\footnotetext{视频画面时间区间：约 18:05--18:45；字幕对应“再乘上另外一个向量叫做 gamma”“再加上 beta”。}

\begin{importantbox}{Gamma 与 beta 的意义}
    BN 不是强迫所有 hidden features 永远均值为 0、方差为 1；它只是先把尺度问题消掉，再把“需要什么尺度和偏移”交还给模型学习。若任务确实需要原来的分布，理论上 $\gamma$ 与 $\beta$ 可以把它学回去。
\end{importantbox}

\subsection{本章小结}

BN 的完整训练时变换是 normalize 后再 affine transform。$\gamma$ 和 $\beta$ 保留表达弹性，让 normalization 不至于成为过强约束。

% =====================================================================
\section{Testing 阶段怎么办}
% =====================================================================

训练时可以用 mini-batch 估计 $\mu_B$ 和 $\sigma_B$。但 testing 或 inference 时，经常一次只来一笔资料。在线服务更不可能等 64 笔 request 凑成 batch 后再统一预测。因此 testing 阶段不能继续依赖当前 batch statistics。

常见解决方法是在 training 时维护 running statistics。每次处理一个 batch 得到 $\mu_t$ 与 $\sigma_t$ 后，更新 moving average：
\[
    \bar\mu \leftarrow p\bar\mu + (1-p)\mu_t,
\]
\[
    \bar\sigma \leftarrow p\bar\sigma + (1-p)\sigma_t.
\]
Testing 时就使用训练过程中累计的 $\bar\mu$ 与 $\bar\sigma$。

\begin{figure}[H]
    \centering
    \includegraphics[width=0.86\textwidth]{figures/fig06_testing_running_average.jpg}
    \caption{Testing 阶段不一定有 batch；实作中通常使用 training 阶段累积的 moving average statistics。\protect\footnotemark}
    \label{fig:testing-running}
\end{figure}
\footnotetext{视频画面时间区间：约 21:50--22:30；字幕对应“training 的时候每一个 batch 计算出来的 mu 跟 sigma 都会拿出来算 moving average”。}

在 PyTorch 中，\texttt{BatchNorm} 层会自动维护类似 \texttt{running\_mean} 与 \texttt{running\_var} 的统计量。调用 \texttt{model.train()} 时，BN 使用当前 batch statistics 并更新 running statistics；调用 \texttt{model.eval()} 时，BN 使用 running statistics，不再用当前 batch 估计。

\begin{warningbox}{常见实作坑}
    验证或测试前忘记调用 \texttt{model.eval()}，会导致 BN 继续使用当前 batch statistics，使结果依赖 batch 组成；训练恢复时忘记切回 \texttt{model.train()}，则 BN 不再更新统计量，也会影响训练行为。
\end{warningbox}

\subsection{本章小结}

BN 在 training 和 testing 阶段行为不同。Training 用当前 batch statistics 并更新 moving average；testing 用训练期间累积的 running statistics。

% =====================================================================
\section{BN 的经验效果与解释争议}
% =====================================================================

原始 Batch Normalization 论文展示了非常直观的实验：加入 BN 后，模型可以更快训练；在某些设定中还可以使用更大的 learning rate。课堂中的图显示，红色和蓝色曲线比未使用 BN 的黑色虚线更快上升。

\begin{figure}[H]
    \centering
    \includegraphics[width=0.86\textwidth]{figures/fig07_bn_original_experiment.jpg}
    \caption{原始 BN 论文中的实验曲线：加入 BN 后训练速度明显加快，并支持更大的 learning rate。\protect\footnotemark}
    \label{fig:bn-experiment}
\end{figure}
\footnotetext{视频画面时间区间：约 23:05--24:35；字幕对应“训练速度显然比黑色虚线快很多”“可以把 learning rate 设大一点”。}

\subsection{Internal Covariate Shift 解释}

BN 最早常被解释为缓解 internal covariate shift。直觉是：前一层参数更新后，后一层看到的 activation distribution 会改变；后一层原本适合旧分布的参数，可能突然不适合新分布。BN 通过让 activation statistics 更稳定，使后一层面对的输入分布变化较小。

\begin{figure}[H]
    \centering
    \includegraphics[width=0.86\textwidth]{figures/fig08_internal_covariate_shift_challenge.jpg}
    \caption{Internal covariate shift 的原始解释：前层更新会改变后层输入分布；但后续实验并不完全支持“BN 主要靠减少分布漂移”这一说法。\protect\footnotemark}
    \label{fig:ics-challenge}
\end{figure}
\footnotetext{视频画面时间区间：约 27:30--28:15；字幕对应“比较训练的时候 A 的分布变化”“不管有没有做 batch normalization，它的变化都不大”。}

后续研究对这个解释提出质疑。实验发现，即使没有明显减少 activation distribution 的变化，BN 仍然能改善训练。因此，“减少 internal covariate shift”可能不是 BN 起作用的全部原因，甚至不一定是主要原因。

\subsection{更可靠的解释：改变 Error Surface}

课堂最后强调的观点是：实验和理论分析更支持 BN 改变了 error surface，让 loss landscape 更平滑、更容易优化。也就是说，BN 的效果可能不是简单让每层输入分布稳定，而是改变了参数空间中 loss 的形状，使 gradient descent 对 learning rate、初始化、局部曲率更不敏感。

\begin{figure}[H]
    \centering
    \includegraphics[width=0.86\textwidth]{figures/fig09_error_surface_serendipity.jpg}
    \caption{后续论文认为 BN 的正面影响可能与 error surface 的改变有关，并称其效果带有 serendipitous 的性质。\protect\footnotemark}
    \label{fig:error-surface}
\end{figure}
\footnotetext{视频画面时间区间：约 28:35--29:20；字幕对应“可以改变 error surface，让 error surface 比较不崎岖”“positive impact 可能是 serendipitous”。}

\begin{knowledgebox}{Serendipitous 的含义}
    老师用盘尼西林的故事解释 serendipity：它不是完全随机的好运，而是在探索中意外发现了有价值的东西。BN 的提出也有类似味道：最初解释未必完全正确，但这个设计在经验上非常有效，并推动了后续 normalization 方法的发展。
\end{knowledgebox}

\subsection{本章小结}

BN 经验上显著改善训练速度和稳定性。早期常用 internal covariate shift 解释它，后续研究更倾向于认为 BN 改变了 loss landscape，使优化更容易。

% =====================================================================
\section{总结与延伸}
% =====================================================================

Batch Normalization 可以从三个层次理解。

第一，作为 feature scale 修正。不同 dimension 的尺度差异会让 error surface 变得狭长崎岖；normalization 把均值和方差拉到相近范围，让 gradient descent 更容易使用统一步长。

第二，作为深层网络内部的动态 normalization。BN 不只处理输入，而是处理 hidden features。训练时，均值与方差来自 mini-batch，因此 BN 层会让 batch 内样本发生统计耦合。

第三，作为优化地形改造。BN 的好处不只是“减少分布漂移”，更可靠的理解是它让 loss landscape 更容易优化，通常支持更大的 learning rate，并提升训练稳定性。

\begin{importantbox}{实作记忆点}
    训练时 BN 用当前 batch 的 $\mu_B,\sigma_B$，并更新 running statistics；测试时 BN 用 running statistics。小 batch 时 BN 可能不稳定；需要时考虑 Layer Normalization、Instance Normalization 或 Group Normalization。
\end{importantbox}

\subsection{延伸阅读}

后续阅读 normalization 方法时，可以比较 Batch Normalization、Layer Normalization、Instance Normalization、Group Normalization 的统计维度：它们的核心差异不是公式长得像不像，而是平均值和方差到底沿哪些 axis 计算。

\end{document}
